%
% 6-argument.tex
%
% (c) 2026 Prof Dr Andreas Müller
%
\section{Das Argumentprinzip
\label{buch:integration:section:argumentprinzip}}
Das Argument einer komplexen Zahl $z$ ist der Winkel $\varphi$,
mit dem $z=|z|\,e^{i\varphi}$ rekonstruiert werden kann.
Dieser ist nur bis auf ein Vielfaches von $2\pi$ bestimmt.
Diese Flexibilität wird benötigt, um die Entwicklung des
Arguments entlang einer Kurve zu untersuchen.
In Abschnitt~\ref{buch:integration:cauchy:subsetion:umlaufzahl}
wurde die Umlaufzahl als Integral mithilfe des Cauchy-Integralsatzes
berechnet.
In diesem Abschnitt soll jetzt gezeigt werden, dass das Bild
einer geschlossenen Kurve die Information darüber enthält, wieviele
Nullstellen und Polstellen eine meromorphe Funktion hat.

\input{chapters/040-analytisch/61-integral.tex}
\input{chapters/040-analytisch/62-meromorph.tex}
\input{chapters/040-analytisch/63-verallgemeinert.tex}
